\section{Limitaciones}

El dump NYIT empleado no incluye tiempos de liberación de tecla: el \textit{dwell time} está ausente y 
el \textit{flight time} coincide con el IKI. La firma no es una rúbrica legal, sino un \textit{proxy} de 
\textit{swipe}. Tecleo y gesto táctil ocurren en ensayos distintos, de modo que no hay fusión temporalmente 
alineada. No hay protocolo de impostor hábil. Un solo modelo de dispositivo (iPhone XR) impide hablar de 
heterogeneidad de hardware. TypeNet, Type2Branch y BehaveFormer se citan como SOTA ajena (Aalto/HuMIdb),
no reentrenada aquí. El único referente del mismo set es Cariello et al.\ (RF/$k$-NN/SVM + mocap);
no hay un Transformer SOTA publicado en NYIT.

La cabeza clásica (Random Forest de rasgos y RF-meta de scores) se evaluó en el mismo leftover y no bate
al Transformer; clonar el protocolo de Cariello con Vicon, lab y ventanas de un segundo sigue fuera de alcance.
El $\tau$ de EER de \textit{dev} del RF de rasgos no se calibra a eval: el ranking se lee por EER, no por ese umbral.

El $N$ de pares genuinos difiere entre el Transformer 3-mod (545, tope de scroll) y el ensemble (2278). 
Comparar EER sin esa nota sobreestima o subestima diferencias. El TTD usa una regla de $k=2$ ventanas 
consecutivas y no modela el tiempo de reloj de una sesión real. El experimento IMU, aunque archivado, no 
replica HuMIdb ni el STDAT de BehaveFormer.

\section{Trabajos futuros}

Las extensiones naturales son: (i)~un corpus con DT/FT reales y, si se mantiene la modalidad de firma, 
trazos de rúbrica en el mismo ensayo que el tecleo; (ii)~impostor hábil y más de un dispositivo; 
(iii)~entrenar y evaluar en Aalto o HuMIdb con el mismo código, sin mezclar ejes con NYIT; (iv)~calibrar 
umbrales por usuario y estudiar TTD en tiempo de pared; (v)~sustituir el \textit{proxy} de firma o 
eliminarlo del sistema propuesto si sigue en azar; (vi)~explotar \textbf{rasgos que el dump ya contiene}
y este pipeline no usa en el sistema propuesto. El \textit{dwell} no se puede inventar: no hay
\textit{release}. Sí se pueden enriquecer los IKI (autocorrelación, histogramas, pausas, FFT), la
cinemática del \textit{swipe} (curvatura, jerk) y, sobre todo, la IMU de teléfono ya extraída como
\textit{sidecar} (tecleo+IMU en leftover: EER $34{,}63\%$, peor que el Transformer 3-mod). En lab hay
además series Vicon y, en algunos trials, \textit{face tracking}; alinear IMU (o mocap) a cada pulsación
como Cariello, sin salir del teléfono, sería la comparación de rasgos pendiente. Sit/walk y casa/lab son
etiquetas de contexto, no identidad nueva.

Recolectar firma legal y bajar el EER mediante reentrenamiento del Transformer quedan fuera del alcance
ya cerrado de esta memoria.
